% 本文件由 paperswarm.tex_gate 从台账自动生成，请勿手工编辑。
% 每个数值均由证据台账注入：claim_id -> (artifact SHA-256, JSON Pointer)。
\section{Introduction}

PaperSwarm is a novel multi-agent system designed to generate reproducible research papers. The system leverages a decentralized architecture where each agent independently contributes to the writing process, ensuring that the final document is both accurate and transparent. To evaluate the efficacy of PaperSwarm, we conducted experiments comparing it against traditional methods, such as ordinary least squares (OLS) and ridge regression.

The experiments were conducted using a dataset with 120 features and 90 training samples per split, with 200 samples reserved for testing. The OLS model, while simple, exhibited significant variability in performance, as indicated by a mean test mean squared error (MSE) of 0.6198 and a standard deviation of 0.1329. In contrast, ridge regression demonstrated greater stability, achieving a mean test MSE of 0.3041 with a much lower standard deviation of 0.0623. This suggests that the selected ridge penalty, 0.1, effectively reduced the variance in model performance.

The relative improvement of ridge regression over OLS was substantial, with a reduction in mean test MSE of 50.94 percent. This improvement is particularly noteworthy as it aligns with the theoretical benefits of ridge regression in mitigating the effects of multicollinearity and overfitting. The irreducible noise floor, 0.1225, serves as a benchmark, highlighting the extent to which the models can reduce error beyond what is inherently present in the data.

These results underscore the importance of considering model complexity and stability in the context of multi-agent systems like PaperSwarm. The high variance observed in OLS and the consistent performance of ridge regression indicate that PaperSwarm can effectively harness the strengths of different methodologies, leading to more robust and reliable research outputs.
